% ============================================================
%  Section 1 — Introduction
% ============================================================

\section{Introduction}
\label{sec:intro}

Fine wine is unusual among investment assets because it ages out. A bottle has a biological
maturity date past which its consumption value declines, yet its market price need not follow
suit: collectors bid for old wine precisely because it can no longer be drunk in the ordinary
sense. This paper uses a large panel of French fine-wine auction transactions to study how
prices evolve over the wine life-cycle, and to show that the shape of the price-age profile
reveals the composition of demand.

\noindent\textit{[Opening paragraph: motivate the puzzle more sharply — most assets appreciate
or depreciate monotonically; wine does neither. Frame the maturity date as the key object:
it creates a natural experiment in which the same bottle transitions from a consumption good
to a collectible in real time.]}

Using 1,068,703 lot-level hammer prices from French fine-wine auctions (1996--2015), we
document a striking non-monotonicity in the price-age profile of premier-quality wines.
Prices rise through the pre-maturity period, peak near the biological maturity window, then
\textit{dip} by approximately 26 percent in a ``trough zone'' spanning the decade after
peak maturity, before recovering to a collector premium at antique ages. This three-phase
pattern is pervasive for Bordeaux Crus Class\'es and Burgundy Premier Crus but is
\textit{absent} for Burgundy Grand Crus, whose auction prices rise monotonically through
the entire life-cycle.

The divergence between Bordeaux and Burgundy Grand Cru identifies the mechanism.
Among young lots (below eight years old), 51.5 percent of Burgundy Grand Cru transactions
exceed \$200 per bottle, compared with 39.0 percent for Bordeaux Crus Class\'es and only
11.2 percent for Burgundy Premier Cru. The high entry price of Burgundy Grand Cru screens
out consumption buyers from the outset: a bidder who intends to drink the wine must be
willing to pay collector-level prices even when the wine is young. There is no consumption
demand segment to exit at peak maturity, and therefore no downward price pressure in the
trough zone.

To formalise this logic, we develop a parsimonious two-type auction model following
\citet{lovo2018}. Buyers are either consumption buyers, who value the wine for drinking
(and whose valuation peaks at maturity and declines thereafter), or collector buyers, who
place an increasing value on rarity and age. The model predicts a non-monotone price-age
profile whenever the consumption-buyer valuation at maturity exceeds the collector-buyer
floor, and a monotone profile when the collector floor is sufficiently high. Three
propositions characterise (i) existence of the non-monotone profile, (ii) the conditions
for trough disappearance, and (iii) comparative statics on trough depth.

Our paper relates to two strands of literature. The first is the empirical literature on
alternative assets as investments \citep{baumol1986, goetzmann1993, mei2002,
korteweg2016, goetzmann2021}. Within this strand, the closest paper is
\citet{dimson2015}, who use a repeat-sales index constructed from Bordeaux Premiers Crus
hammer prices over 1900--2012 to estimate a net real return to wine of 4.1\,\% per annum
and document that young maturing wines generate the highest returns while past-maturity
ch\^{a}teaux provide growing non-pecuniary (collector) benefits. Our paper departs from
\citet{dimson2015} in three important ways. First, they use raw calendar age and therefore
conflate wines at very different stages of their biological life-cycle: a twenty-year
Bordeaux Premier Cru may be at peak maturity while a twenty-year village wine is already in
the trough zone. We normalise age by each wine's individual maturity date, obtained from
Tastingbook for 75.9\,\% of lots and imputed otherwise, which maps heterogeneous calendar
ages onto a common life-cycle grid and isolates the maturity transition as the key event.
Second, the repeat-sales method used in \citet{dimson2015} requires a bottle to be resold
at least twice and therefore systematically excludes the most illiquid segments of the
market; our lot-level panel is ten times larger and covers three wine regions. Third, and
most importantly, \citet{dimson2015} characterise the maturity kink descriptively but do not
model the mechanism. The decomposition of the price dip into a buyer composition effect ---
the exit of consumption buyers at peak maturity --- and the ensuing recovery as collectors
dominate antique-age auctions requires a structural model with heterogeneous buyer types,
which we provide. \citet{breeden2017} use age-period-cohort decomposition on 1.5 million
auction records and confirm the three-phase price profile, but likewise stop short of
modelling why the dip is present for some wine categories and absent for others.

The second strand is the theory of heterogeneous-buyer markets for durable goods. The
structural model of \citet{lovo2018} is our closest theoretical antecedent. They develop
a two-type framework for the art market in which an asset is held by either a ``pleasure
buyer'' (who derives a flow utility from ownership) or a ``speculative buyer'' (who holds
it purely for resale). Their key result is that speculation generates price cycles and
excess volatility relative to fundamental value. Our model shares the two-type architecture
but differs in a central respect: \textit{wine has a finite consumption life}. We model the
consumption-buyer valuation as single-peaked at the maturity date and declining thereafter
--- a restriction that has no counterpart in \citet{lovo2018}, where the pleasure buyer's
flow utility is constant and the asset never expires. It is precisely this biological
maturity structure that generates the non-monotone price-age profile and the composition
effect at the heart of our analysis. A second difference is institutional: wine is sold in
one-shot Vickrey auctions rather than in a search-and-matching market, which allows us to
derive the equilibrium price in closed form and take it directly to data. The maturity
date also provides a natural source of cross-sectional variation to identify the model: the
same wine category at the same calendar age can be pre-maturity or post-maturity depending
on the wine's individual $\tau_i$, and this variation is observable.

\noindent\textit{[Roadmap paragraph: Section~\ref{sec:data} describes the data.
Section~\ref{sec:empirics} presents the empirical strategy. Section~\ref{sec:results} reports
results. Section~\ref{sec:model} develops the structural model that rationalises these
findings. Section~\ref{sec:conclusion} concludes.]}
